% ===================================================================
%  第 6 节 权重 3/2 处的月光（Moonshine at weight three-halves）
%  原文：TASI10.tex 第 2073--2124 行（PDF p.55--57）
% ===================================================================
\section{权重 $3/2$ 处的月光}
\label{ch:moonshineweight32}

\begin{keytakeaways}
  \item 权重 $3/2$ 的月光目前只有两个确立的例子：\emph{O'Nan 月光}把一组权 $3/2$ 模拟模形式与 O'Nan 群相联系；\emph{类数月光}（class number moonshine）则由 Hurwitz–Kronecker 类数构造出容纳无穷多个群 $\ZZ_p$ 的一族模拟模形式。
  \item 与权 $0$、权 $1/2$ 的月光不同，权 $3/2$ 的模拟模形式不由模性与解析性唯一确定，其 Fourier 系数显著地受权 $3/2$ 尖形式支配。
  \item 超模在这里是内在而不可避免的：同一 $q$ 幂的 Fourier 系数可以是群表示迹的差，因而必须使用超迹。
  \item 权 $3/2$ 尖形式通过 Shimura 对应与 Waldspurger 型关系连接椭圆曲线的特殊 $L$ 值，把月光与算术几何、Tate–Shafarevich 群联系起来。
  \item 月光型联系可能普遍存在：除 pariah 散在群之外，有限单群分类中非散在单群的一个无限族 $\ZZ_p$ 也参与其中；这给物理提出了构造相应超模的挑战。
\end{keytakeaways}

这是一个当前发展迅速的月光研究领域。因此我们将写得非常简略，只突出近期发展的若干关键特征与前景。

% [Tutor Note: 权重 $k$ 的模形式在 Serre 对偶下的配对权重是 $2-k$，所以 $w=1/2$ 的对偶权重为 $2-1/2=3/2$。原文写成 $w=2-\frac{1}{2}=\frac{3}{2}$ 正是强调这一点；不要误记为 $w=1/2+1$，也不要把 $2-w$ 与 $1-w$ 混用。]

在对权重 $1/2$ 的月光作了富有成果的探索之后，人们自然会追问对偶权重——
$w=2-\frac{1}{2}=\frac{3}{2}$——的月光图景。
权重 $3\over 2$ 的模拟模形式所受的约束远少于权重 $1\over 2$ 的对应物，
结果它们为研究有限群联系提供了肥沃的土壤，已经带给我们
\emph{O'Nan 月光} \cite{onan} 与 \emph{类数月光} \cite{class_nr_moonshine}；
后者涉及 \eq{hurwitz_mock} 中描述的 Hurwitz–Kronecker 类数模拟模形式。

% [Tutor Note: Hurwitz–Kronecker 类数 $H(N)$ 是判别式 $-N$ 的整二次型在 $\PSL_2(\ZZ)$ 下的类数（按自同构群大小加权），其生成函数 ${\cal H}(\tau)=\sum_{N\ge 0}H(N)q^N$ 是权 $3/2$、水平 $4$ 的模拟模形式，影子正比于 theta 函数 $\sum_{n\in\ZZ}q^{n^2}$。它就是 \eq{hurwitz_mock} 中的类数模拟模形式；注意它不等于普通类数 $h(-N)$，两者相差一个权因子。]

权重 $3/2$ 中多半还有更多宝藏。这里我们只评论两个已经确立的月光例子。
我们将指出权重 $3/2$ 的这两个月光例子所具有的三个显著性质\footnote{事实上，这些性质使权重 $3/2$ 的联系与迄今讨论的权重零与权重 $1/2$ 的月光例子如此不同，以至于人们会怀疑是否还应当对它们使用“月光”一词。}。

% [Tutor Note: 作者说“多半还有更多宝藏”而未给论证。实质原因是权 $1/2$ 的模拟模形式受 Serre–Stark 型定理的强约束（全纯部分几乎由 theta 函数决定），而权 $3/2$ 没有相应定理，可承载的群表示自由度大得多；因此“更多例子存在”是一个经验判断，不是定理。]

\myparagraph{超模}

在 \S\ref{cha:weight0} 讨论的魔群月光与 Conway 月光中，底层的模是“实”模（相对于超模 \eq{def:supermodule}），
所涉及的群特征标是普通迹（相对于超迹）。
Mathieu 月光与本影月光中也是如此，只是 \eq{rep_mock_umbral} 中的极点项例外：
它们的系数为负，却不携带任何关于群表示的非平凡信息。
在 Thompson 月光中我们遇到超模，但方式相当平凡：超模是纯偶的还是纯奇的，
取决于相应的 $q$ 幂在模 $4$ 下是 $+1$ 还是 $-1$。
特别地，在权重 $1/2$ 中我们不会遇到这样的月光例子：模拟模形式的 Fourier 系数
由群表示迹的差给出，也就是说，在每个齐次分量中超模总是要么为偶、要么为奇。

% [Tutor Note: 超模 \eq{def:supermodule} 指 $\ZZ_2$ 分次模，偶分量的维数减奇分量的维数给出超迹 Str。在权 $1/2$ 的 Thompson 月光中，按 $q$ 幂模 $4$ 的取值 $+1$ 或 $-1$，每个齐次分量整体取偶或整体取奇，所以普通迹已经够用；权 $3/2$ 中同一齐次分量可以同时含偶、奇部分，普通迹会丢掉信息，必须使用超迹。]

在权重 $3/2$ 中这不再成立：超模的“超”性质是我们讨论的这两个月光例子内在且不可避免的特征。

\myparagraph{尖形式}

在迄今讨论的月光例子中，所涉及的（模拟）模对象都十分特殊：它们本质上由自身的模性质与解析性质唯一确定。
特别地，魔群月光与 Conway 月光的分次特征标（graded characters）是 Hauptmodul，
在一个可加常数之内可以写成只含单个极点项 $q^{-1}$ 的 Rademacher 和。
本影月光（包括 Mathieu 月光）与 Thompson 月光的分次特征标，在加上 theta 函数之后也是（投影）Rademacher 和
（参见 \eq{thompson_mckaythompson} 与 \cite{cheng2018weight}）。
后者可以说只是一个小修正：除极少数指数外，theta 函数的系数都为零
（除非指数在某个特定的重标度下为平方数），而且它们不携带任何有趣的数论信息。

% [Tutor Note: “由模性与解析性本质唯一确定”一语省略了论证：给定权、给定同余子群与给定极点结构（此处是单个 $q^{-1}$），Rademacher 和给出唯一的模对象；极点结构本身又由可复制性（replicability）约束固定。这正是权 $0$ 与权 $1/2$ 月光能用少量数据确定整条 $q$ 级数的原因。]

% [Tutor Note: theta 函数的系数只在（重标度后的）平方数处非零，所以“加上 theta 函数”只修正极少数 $q$ 幂。但 theta 函数系数本身不含群表示信息，不要把它当成月光内容的一部分，也不要据此认为整个模对象仍未确定。]

这个关键的唯一性性质曾用来区分权重零与权重 $1/2$ 月光的模拟模形式；
在权重 $3/2$ 的月光中，它不再是相关的判据。
O'Nan 月光与 类数月光中的模拟模形式，以非常显著的方式涉及权重 $3/2$ 的尖形式。
（尖形式的定义见 \S\ref{subsec:modular_forms} 末尾。）

这一性质打开了与算术几何——数学中一个完全不同的领域——相联系的大门。
这一联系的主要要素是：1）在所谓的 Shimura 对应 \cite{MR783554} 之下，
权重 $3/2$ 新形式的系数与相应权重二新形式的特殊 $L$ 值之间的 Waldspurger 型关系。
2）把椭圆曲线的 $L$ 函数与某些权重二新形式的 $L$ 函数联系起来的模性定理。
3）Birch–Swinnerton-Dyer 猜想（的强形式）（及其已证明的部分）：
它把椭圆曲线的特殊 $L$ 值与它们的重要算术数据——包括 Tate–Shafarevich 群的阶——联系起来。
参见 \cite{MR833192,MR954295,MR3148103}。
4）月光模对象之间的同余，源自特征标之间的同余，后者是有限群表示的基本性质。

% [Tutor Note: Waldspurger 公式给出的是系数平方与特殊 $L$ 值的比例关系（通常形如 $|a_f(|D|)|^2$ 正比于 $L(f,1)$，差 $|D|$ 的幂次与局部因子），并非系数本身与 $L$ 值成线性正比；此外权 $3/2$ 新形式与权 $2$ 新形式之间的对应是 Shimura 对应，不是恒等映射，引用 \cite{MR783554} 时勿把方向弄反。]

% [Tutor Note: 作者把 BSD 列为第三项，是因为 BSD 把 $L(E,1)$ 的零点阶与调节子、挠子群、以及 Tate–Shafarevich 群的阶联系起来；与第 1、2 项合用，才能把权 $3/2$ 系数之间的同余翻译成该群 $p$ 部分的陈述。注意这里同时用到 BSD 的“强形式”与“已证部分”，结论中凡依赖未证部分者都是条件性的。]

综合以上要素，权重 $3/2$ 尖形式与有限群之间的联系引出一系列无穷多的、
关于 Tate–Shafarevich 群 $p$ 部分的陈述，而这些陈述由有限群结构组织起来。
此外，在 类数月光中，附着于某些模数据的椭圆曲线（以及一般的其他阿贝尔簇）的挠子群，
在构造月光模拟模形式时起重要作用。
探索模拟模形式中的有限群结构在算术几何中的含义，是月光研究中一个有趣的新数学方向。

\myparagraph{新的有限群}

上述两个性质分别涉及月光联系的有限群一侧与模形式一侧，
它们使权重 $3/2$ 月光有可能把新类型的有限群纳入月光的领域。

O'Nan 月光，顾名思义，把一组权重 $3/2$ 的模拟模形式与 O'Nan 群联系起来。
O'Nan 群于 1976 年在有限群分类的紧张发展中被发现，是与魔群无关的 6 个 pariah 群之一。
另一方面，类数月光给出由 Hurwitz–Kronecker 类数 \eq{def:hurwitz} 产生的
无穷多个权重 $3/2$ 模拟模形式构成的一族，它容纳\emph{无穷多个群}的表示。
回忆一下：在有限单群的分类中，除散在群之外还有三个无限族，
其中之一是素数阶 $p$ 的循环群 $\ZZ_p$。

% [Tutor Note: 原文 “cycle groups” 应为 “cyclic groups”（循环群）之误：有限单群分类中的三个无限族是 $\ZZ_p$（$p$ 为素数）、交错群 $A_n$（$n\ge 5$）与 Lie 型（Chevalley）群；本句随后只用到 $\ZZ_p$。译文按“循环群”处理，原词登记为疑似笔误，页面图像 PDF p.56 亦作 cycle groups。]

% [Tutor Note: 6 个 pariah 群指不参与魔群的散在单群：$J_1$、$J_3$、$J_4$、$\mathrm{Ly}$、$\mathrm{O'N}$、$\mathrm{Ru}$。O'Nan 群由 O'Nan 于 1976 年构造，阶为 $2^9\cdot 3^4\cdot 5\cdot 7^3\cdot 11\cdot 19\cdot 31$。它出现在魔群月光之外，正是本节强调“月光不限于魔群”的关键例证。]

结果表明，对每个素数 $p$，都有一个水平 $4$ 的模拟模形式 ${\cal H}_{p,1}$——
它是 类数模拟模形式 ${\cal H}$ \eq{hurwitz_mock} 的标量倍——
以及一个水平 $4p$ 的模拟模形式 ${\cal H}_{p,p}$，
使得它们与 $\ZZ_p$ 的一个 $\ZZ$ 分次超模的两个分次特征标相符。

% [Tutor Note: 对 $\ZZ_p$ 而言不可约表示都是一维的（特征标取 $p$ 次单位根），所以两个分次特征标足以编码整个超模；水平 $4$ 与 $4p$ 正对应未扭与 $\ZZ_p$ 扭结扇区（twisted sector）的模参数，因此 $\ZZ_p$ 的“无穷族”恰好对应一族按 $p$ 增长的模形式。]

除了这个无限的群族之外，\cite{class_nr_moonshine} 还表明，
不同水平（由群元阶的 4 倍给出）的模拟模形式共同编码了有趣的 Mathieu 群的超模。

这些权重 $3/2$ 月光的实例大大拓宽了模对象与有限群之间联系的视野。
权重零与权重 $1/2$ 的所有月光实例都与和魔群有关的群相联系，
但 O'Nan 例子与类数例子表明，情况未必如此。
在这些例子中已经证明：一个 pariah 散在群，以及三个非散在有限单群无限族之一，
也以类似方式与（模拟）模形式相联系！
这暗示了一种诱人的可能性：月光型联系其实是普遍的，是模对象与有限群图景中的内在性质；
而有限群是模对象的一种重要组织性结构。

最后但同样重要的是，这些涉及无穷多个群的超模的新联系给物理提出了有趣的挑战。
物理还能提供这些有限群模并解释这些联系吗？
如果能够，哪些相关物理系统可以产生这些超模？
